%%%%%%%%%%%%%%%%%%%%%%%%%%%%%%%%%%%%%%%%%%%%%%%%%%%%%%%%%%%%%%%%%%%%%%%%%%%%%%%%
% Document LaTeX : Équations différentielles - Version simplifiée
% Auteur : Laurent
% Date : 2026-07-11
% Description : Exemples basiques d'équations différentielles
%%%%%%%%%%%%%%%%%%%%%%%%%%%%%%%%%%%%%%%%%%%%%%%%%%%%%%%%%%%%%%%%%%%%%%%%%%%%%%%

\documentclass[11pt]{article}
\usepackage[utf8]{inputenc}
\usepackage[french]{babel}
\usepackage{amsmath}
\usepackage{amssymb}
\usepackage{geometry}

% Configuration de la page
\geometry{a4paper, margin=1.5cm}

% Titre et métadonnées
\title{Exemples Simples d'Équations Différentielles}
\author{Laurent}
\date{\today}

%%%%%%%%%%%%%%%%%%%%%%%%%%%%%%%%%%%%%%%%%%%%%%%%%%%%%%%%%%%%%%%%%%%%%%%%%%%%%%%
\begin{document}

\maketitle

\section{Équations du 1er ordre}

\subsection{Variables séparables}

\textbf{Exemple 1 :} Résoudre $y' = 2xy$.

\textbf{Solution :}
\begin{align*}
    \frac{dy}{dx} &= 2xy \\
    \frac{dy}{y} &= 2x \, dx \\
    \int \frac{1}{y} \, dy &= \int 2x \, dx \\
    \ln|y| &= x^2 + C \\
    y &= \pm e^{x^2 + C} = C' e^{x^2}
\end{align*}

\textbf{Réponse finale :} $\boxed{y = C e^{x^2}}$

\bigskip

\textbf{Exemple 2 :} Résoudre $y' = \frac{y}{x}$ avec $y(1) = 2$.

\textbf{Solution :}
\begin{align*}
    \frac{dy}{dx} &= \frac{y}{x} \\
    \frac{dy}{y} &= \frac{dx}{x} \\
    \int \frac{1}{y} \, dy &= \int \frac{1}{x} \, dx \\
    \ln|y| &= \ln|x| + C \\
    y &= \pm e^C x = Cx
\end{align*}
Avec $y(1) = 2$ : $2 = C \cdot 1 \implies C = 2$.

\textbf{Réponse finale :} $\boxed{y = 2x}$

\newpage

\subsection{Linéaires du 1er ordre}

\textbf{Exemple 3 :} Résoudre $y' + 2y = 4$.

\textbf{Solution :}
Facteur intégrant : $\mu(x) = e^{\int 2 \, dx} = e^{2x}$.
\begin{align*}
    e^{2x} y' + 2 e^{2x} y &= 4 e^{2x} \\
    \frac{d}{dx}(e^{2x} y) &= 4 e^{2x} \\
    e^{2x} y &= 2 e^{2x} + C \\
    y &= 2 + C e^{-2x}
\end{align*}

\textbf{Réponse finale :} $\boxed{y = 2 + C e^{-2x}}$

\bigskip

\textbf{Exemple 4 :} Résoudre $y' - 3y = e^{2x}$.

\textbf{Solution :}
Facteur intégrant : $\mu(x) = e^{\int -3 \, dx} = e^{-3x}$.
\begin{align*}
    e^{-3x} y' - 3 e^{-3x} y &= e^{-x} \\
    \frac{d}{dx}(e^{-3x} y) &= e^{-x} \\
    e^{-3x} y &= -e^{-x} + C \\
    y &= -e^{2x} + C e^{3x}
\end{align*}

\textbf{Réponse finale :} $\boxed{y = -e^{2x} + C e^{3x}}$

\newpage

\section{Équations du 2nd ordre}

\subsection{À coefficients constants}

\textbf{Exemple 5 :} Résoudre $y'' - 3y' + 2y = 0$.

\textbf{Solution :}
Équation caractéristique : $r^2 - 3r + 2 = 0$.
\begin{align*}
    r &= \frac{3 \pm \sqrt{9 - 8}}{2} = \frac{3 \pm 1}{2} \\
    r_1 &= 2, \quad r_2 = 1
\end{align*}

\textbf{Réponse finale :} $\boxed{y = C_1 e^{2x} + C_2 e^{x}}$

\bigskip

\textbf{Exemple 6 :} Résoudre $y'' + 4y = 0$.

\textbf{Solution :}
Équation caractéristique : $r^2 + 4 = 0 \implies r = \pm 2i$.

\textbf{Réponse finale :} $\boxed{y = C_1 \cos(2x) + C_2 \sin(2x)}$

\bigskip

\textbf{Exemple 7 :} Résoudre $y'' - 4y' + 4y = 0$.

\textbf{Solution :}
Équation caractéristique : $r^2 - 4r + 4 = 0 \implies (r - 2)^2 = 0 \implies r = 2$ (racine double).

\textbf{Réponse finale :} $\boxed{y = (C_1 + C_2 x) e^{2x}}$

\newpage

\subsection{Non homogènes}

\textbf{Exemple 8 :} Résoudre $y'' + y = \sin(x)$.

\textbf{Solution :}
Solution homogène : $y_h = C_1 \cos(x) + C_2 \sin(x)$.

On cherche une solution particulière de la forme $y_p = x (A \cos(x) + B \sin(x))$ (car $\sin(x)$ est solution de l'homogène).
\begin{align*}
    y_p &= x (A \cos(x) + B \sin(x)) \\
    y_p' &= A \cos(x) + B \sin(x) + x (-A \sin(x) + B \cos(x)) \\
    y_p'' &= -A \sin(x) + B \cos(x) - A \sin(x) + B \cos(x) + x (-A \cos(x) - B \sin(x)) \\
    y_p'' + y_p &= -2A \sin(x) + 2B \cos(x) = \sin(x)
\end{align*}
On identifie : $-2A = 1 \implies A = -\frac{1}{2}$ et $2B = 0 \implies B = 0$.

Donc $y_p = -\frac{x}{2} \cos(x)$.

\textbf{Réponse finale :} $\boxed{y = C_1 \cos(x) + C_2 \sin(x) - \frac{x}{2} \cos(x)}$

\newpage

\section{Systèmes d'équations différentielles}

\textbf{Exemple 9 :} Résoudre le système :
\begin{align*}
    x' &= x + y \\
    y' &= -x + y
\end{align*}

\textbf{Solution :}
On écrit sous forme matricielle :
\[
\begin{pmatrix} x \\ y \end{pmatrix}' = \begin{pmatrix} 1 & 1 \\ -1 & 1 \end{pmatrix} \begin{pmatrix} x \\ y \end{pmatrix}
\]
Valeurs propres : $\det \begin{pmatrix} 1 - \lambda & 1 \\ -1 & 1 - \lambda \end{pmatrix} = (1 - \lambda)^2 + 1 = \lambda^2 - 2\lambda + 2 = 0$.
\begin{align*}
    \lambda &= \frac{2 \pm \sqrt{4 - 8}}{2} = 1 \pm i
\end{align*}
Vecteurs propres pour $\lambda = 1 + i$ :
\[
\begin{pmatrix} -i & 1 \\ -1 & -i \end{pmatrix} \begin{pmatrix} v_1 \\ v_2 \end{pmatrix} = 0 \implies v_1 = i, \quad v_2 = 1
\]
Solution générale :
\[
\begin{pmatrix} x \\ y \end{pmatrix} = e^{t} \left( C_1 \begin{pmatrix} \cos(t) \\ -\sin(t) \end{pmatrix} + C_2 \begin{pmatrix} \sin(t) \\ \cos(t) \end{pmatrix} \right)
\]

\textbf{Réponse finale :}
\[
\boxed{\begin{cases}
    x = e^{t} (C_1 \cos(t) + C_2 \sin(t)) \\
    y = e^{t} (-C_1 \sin(t) + C_2 \cos(t))
\end{cases}}
\]

\end{document}
